% Rainbow Run — Analytics Event Playbook (Dub + PostHog), scenario by scenario.
% TeX backend (Tectonic). Build: node tools/pdf/tex-build.mjs docs/pdf/rr-event-playbook.tex
\documentclass[11pt]{article}

\usepackage[a4paper,margin=0.85in]{geometry}
\usepackage{graphicx}
\usepackage{xcolor}
\usepackage{booktabs}
\usepackage{array}
\usepackage{enumitem}
\usepackage{titlesec}
\usepackage{fancyhdr}
\usepackage{listings}
\usepackage{microtype}
\usepackage{tikz}
\usetikzlibrary{arrows.meta}
\usepackage[colorlinks=true,linkcolor=accent,urlcolor=accent]{hyperref}

\newcommand{\doctitle}{Rainbow Run --- Analytics Event Playbook}

\definecolor{ink}{HTML}{1B1F24}
\definecolor{accent}{HTML}{0B7285}
\definecolor{dub}{HTML}{E8590C}
\definecolor{ph}{HTML}{6741D9}
\definecolor{hub}{HTML}{2B8A3E}
\definecolor{soft}{HTML}{F1F3F5}
\definecolor{line}{HTML}{CED4DA}

\titleformat{\section}{\Large\bfseries\color{accent}}{\thesection}{0.6em}{}
\titleformat{\subsection}{\large\bfseries\color{ink}}{}{0em}{}
\titlespacing*{\section}{0pt}{1.2em}{0.4em}
\titlespacing*{\subsection}{0pt}{0.9em}{0.3em}

\pagestyle{fancy}\fancyhf{}
\renewcommand{\headrulewidth}{0pt}\renewcommand{\footrulewidth}{0.4pt}
\fancyfoot[L]{\small\color{ink}Rainbow Run · event playbook}
\fancyfoot[R]{\small\thepage}

\newcommand{\callout}[1]{\par\smallskip\noindent
  \colorbox{soft}{\parbox{\dimexpr\linewidth-2\fboxsep\relax}{\color{ink}#1}}\par\smallskip}
\newcolumntype{P}[1]{>{\raggedright\arraybackslash}p{#1}}
\setlength{\emergencystretch}{3em}\tolerance=1200

% layer chips
\newcommand{\Dub}{\textcolor{dub}{\textbf{Dub}}}
\newcommand{\PH}{\textcolor{ph}{\textbf{PostHog}}}
\newcommand{\Hub}{\textcolor{hub}{\textbf{Hub}}}

\lstdefinestyle{json}{basicstyle=\ttfamily\scriptsize,backgroundcolor=\color{soft},
  frame=single,rulecolor=\color{line},breaklines=true,columns=fullflexible,
  xleftmargin=0.5em,framexleftmargin=0.5em,aboveskip=4pt,belowskip=4pt}

\title{\color{accent}\doctitle\\[2pt]\large\color{ink}\normalfont What \textcolor{dub}{Dub} \& \textcolor{ph}{PostHog} record, scenario by scenario}
\author{Engine · TeX backend}
\date{\today}

\begin{document}
\maketitle\thispagestyle{fancy}

\section{How to read this}
Every player interaction leaves a trail across \emph{three} layers. This playbook walks real
Rainbow Run journeys and shows exactly what each layer captures at each step --- so you know what
data you'll have to answer ``what's working?''.

\begin{center}\small
\begin{tabular}{@{}P{2.4cm}P{5.0cm}P{7.1cm}@{}}
\toprule
Layer & Records & Answers \\
\midrule
\Dub & the \textbf{click} on a short's link --- device, geo, referrer, \texttt{dub\_id} & which \emph{short} drives traffic; (phase 2) click$\to$conversion \\
\PH & the \textbf{in-game funnel} --- every event below, with \texttt{utm\_*}/\texttt{dub\_id} attached, + session replay & where players drop, who converts, who returns \\
\Hub \texttt{/api/track} & a \textbf{server-side count} of the same events (ad-blocker-proof) & the true denominator; the game-design eval signal \\
\bottomrule
\end{tabular}
\end{center}

\callout{\textbf{The thread that ties it together:} a short's Dub link lands on
\texttt{/go/rainbow-run?utm\_content=L1-sunbeam-meadow} (Dub appends \texttt{dub\_id}); the game
reads those off the URL and stamps them onto \emph{every} \PH{} event and \Hub{} payload. So every
chart can be sliced by \textbf{which short the player came from}.}

\section{The event vocabulary (grounded in the real game)}
Rainbow Run is 5 weather biomes --- \emph{Sunbeam Meadow, Drizzle Dunes, Rainbow Bridge,
Frostspark Peaks, Thunder Heights} --- played by \emph{Labooboo \& Siya}. The trackable signals:

\begin{center}\footnotesize\setlength{\tabcolsep}{4pt}
\begin{tabular}{@{}P{2.7cm}P{4.6cm}P{6.7cm}@{}}
\toprule
Event & Fires when & Properties (the interesting ones) \\
\midrule
\texttt{game\_open} & first load & utm\_source/medium/\textbf{content}, dub\_id, \$device\_type \\
\texttt{level\_start} & a biome (re)loads & level (1--5), \textbf{biome} \\
\texttt{coin\_collect} & a coin is grabbed & level, coins\_so\_far, x, y \\
\texttt{weather\_shift} & biome weather cycles (sun$\to$rainbow, $\to$storm \ldots) & level, weather \\
\texttt{player\_death} & a death & level, biome, \textbf{cause} (fall/lagoon/enemy/crumble/ice), x \\
\texttt{level\_complete} & the goal is reached & level, duration\_ms, deaths, coins \\
\texttt{game\_complete} & Thunder Heights cleared & total\_time\_ms, total\_deaths \\
\texttt{player\_idle} / \texttt{player\_quit} & no input $>$Ns / tab closed & level, x, ms \\
\bottomrule
\end{tabular}
\end{center}

\section{Scenarios --- how it actually plays out}

% ---------- Scenario 1 ----------
\subsection{1 · The full conversion (the happy path)}
A player on their phone taps the \emph{Thunder Heights} YouTube Short, plays from the start, and
finishes all five biomes.

\newcommand{\step}[2]{\item[\textcolor{#1}{$\bullet$}]}
\begin{enumerate}[leftmargin=2.2em,itemsep=1pt,topsep=2pt]
  \item[\textcolor{dub}{1.}] taps \texttt{dub.sh/PPDuC2r} $\to$ \Dub{} logs a \textbf{click} (device=Mobile, geo, ref=youtube, mints \texttt{dub\_id})
  \item[\textcolor{hub}{2.}] redirect hits \texttt{/go/rainbow-run?utm\_content=L5-thunder-heights} $\to$ \Hub{} logs \texttt{click rainbow-run|youtube}
  \item[\textcolor{ph}{3.}] game loads $\to$ \PH{} \texttt{\$pageview}, then \texttt{game\_open} (carrying utm + dub\_id + device)
  \item[\textcolor{ph}{4.}] \texttt{level\_start} L1 Sunbeam Meadow $\to$ 7$\times$ \texttt{coin\_collect} $\to$ \texttt{weather\_shift} (sun$\to$rainbow) $\to$ \texttt{level\_complete} (0 deaths)
  \item[\textcolor{ph}{5.}] L2$\to$L4 \ldots one \texttt{player\_death} (cause=\texttt{crumble}) on Frostspark Peaks, respawns, \texttt{level\_complete}
  \item[\textcolor{ph}{6.}] clears L5 $\to$ \texttt{game\_complete} (total\_deaths=1). \Hub{} mirrors every event server-side.
\end{enumerate}
\noindent What \PH{} stores for step 3 (every later event inherits the registered attribution):
\begin{lstlisting}[style=json]
{ "event": "game_open",
  "properties": { "utm_source":"youtube", "utm_medium":"short",
    "utm_content":"L5-thunder-heights", "dub_id":"pPDuC2r_x9",
    "$device_type":"Mobile", "$current_url":".../?utm_content=L5-thunder-heights" } }
\end{lstlisting}
\noindent and the matching \Hub{} call (ad-blocker-proof, same event, server-side):
\begin{lstlisting}[style=json]
POST /api/track  {"game":"rainbow-run","event":"game_open","source":"youtube"}
\end{lstlisting}
\callout{\textbf{This player counts as a full conversion under \texttt{utm\_content=L5}} --- they
land in the funnel's bottom step \emph{and} in ``completion rate by short'' for the Thunder
Heights short.}

% ---------- Scenario 2 ----------
\subsection{2 · The top-of-funnel bounce (click, no play)}
Taps the short, the page opens, but they close the tab before moving.
\begin{itemize}[leftmargin=2.2em,itemsep=1pt,topsep=2pt]
  \item \Dub{} \textbf{click} ✓, \Hub{} \texttt{click} ✓ --- but \PH{} sees only \texttt{\$pageview} + maybe \texttt{game\_open}, then \texttt{\$pageleave}. No \texttt{level\_start}.
\end{itemize}
\textbf{What it tells you:} the funnel drops hard at \texttt{game\_open $\to$ level\_start}. If one
short bounces far worse than others (slice by \texttt{utm\_content}), the \emph{landing} is the
problem --- slow load, wrong expectation set by the short --- not the game.

% ---------- Scenario 3 ----------
\subsection{3 · The struggler (a difficulty signal)}
Comes from the \emph{Drizzle Dunes} short, hits the ice + conveyor stretch, dies four times, quits.
\begin{lstlisting}[style=json]
player_death { level:2, biome:"Drizzle Dunes", cause:"ice",   x:520 }
player_death { level:2, biome:"Drizzle Dunes", cause:"fall",  x:540 }   // x4, clustered at the belt
player_idle  { level:2, x:300, ms:9000 }
player_quit  { level:2, x:300 }                                          // no level_complete for L2
\end{lstlisting}
\textbf{What it tells you:} \texttt{deaths-by-level} shows a spike on L2 at \texttt{x$\approx$520}
(the ice belt). Open the \PH{} \textbf{session replay} of that play to watch it; the death cluster
is the difficulty heatmap. This is the loop back into game design --- widen the belt, slow the
conveyor --- and the same event feeds the engine's reach/difficulty eval.

% ---------- Scenario 4 ----------
\subsection{4 · The returner (retention by short)}
Day 1: arrives from the \emph{Sunbeam Meadow} short, plays L1--L2, leaves. Day 3: opens the game
directly (\texttt{utm\_source=direct}) and plays L3.
\begin{itemize}[leftmargin=2.2em,itemsep=1pt,topsep=2pt]
  \item \PH{} ties both visits to one \texttt{distinct\_id}; the day-1 \texttt{utm\_content} stays on the person as a property.
  \item \PH{} \textbf{Retention} shows a Day-3 return, attributable to the \emph{Sunbeam Meadow} short.
\end{itemize}
\textbf{What it tells you:} which shorts bring \emph{sticky} players, not just clicks. A short with
fewer but higher-retention plays can beat a viral-but-bouncy one.

% ---------- Scenario 5 ----------
\subsection{5 · Mobile vs desktop (a device gap)}
A TikTok short sends mostly mobile traffic; those players drop at \texttt{level\_start} far more
than desktop.
\begin{center}\footnotesize
\begin{tabular}{@{}lrrr@{}}
\toprule
device & game\_open & level\_complete & completion \% \\
\midrule
Desktop & 120 & 54 & 45.0 \\
Mobile  & 300 & 39 & 13.0 \\
\bottomrule
\end{tabular}
\end{center}
\textbf{What it tells you:} \texttt{convert-by-device} exposes a touch-controls problem. Fixes:
tune the on-screen controls, or use \Dub{} \emph{device targeting} to route mobile to a
touch-optimized entry.

% ---------- Scenario 6 ----------
\subsection{6 · The ad-blocked player (why the Hub backstop exists)}
A player on uBlock: the \PH{} \texttt{phc\_} ingestion request is blocked, so \emph{zero} client
events arrive --- but the same-origin \Hub{} \texttt{/api/track} calls still land.
\begin{center}\footnotesize
\begin{tabular}{@{}lrr@{}}
\toprule
 & \PH{} (client) & \Hub{} (server) \\
\midrule
plays counted & 70 & 100 \\
\bottomrule
\end{tabular}
\end{center}
\textbf{What it tells you:} reconciling the two sizes your blind spot (here $\approx$30\% of
short-form, mobile-heavy traffic). \PH{} is your rich funnel; \Hub{} is the honest denominator.

\section{What you can now answer (scenario $\to$ insight $\to$ query)}
\begin{center}\footnotesize\setlength{\tabcolsep}{4pt}
\begin{tabular}{@{}P{3.0cm}P{5.2cm}P{5.8cm}@{}}
\toprule
Question & PostHog surface & HogQL (\texttt{scripts/posthog-query.mjs}) \\
\midrule
Which short converts best? & Funnel, breakdown \texttt{utm\_content} & \texttt{--canned funnel-by-short} \\
Where do players die? & Trends/Paths + session replay & \texttt{--canned deaths-by-level} \\
Mobile vs desktop? & Funnel, breakdown \$device\_type & \texttt{--canned convert-by-device} \\
Which short retains? & Retention, breakdown \texttt{utm\_content} & persons + first-touch query \\
Who converts (clusters)? & Funnel \textbf{correlation} + behavioral cohorts & \texttt{GROUP BY} on segments \\
Ad-block blind spot? & --- (compare to Hub) & hub \texttt{/api/funnel} vs PH counts \\
\bottomrule
\end{tabular}
\end{center}

\callout{\textbf{Status:} the read path is live (the \texttt{posthog-query.mjs} tool + your personal
key already query project 484401). The events above start flowing the moment Rainbow Run is
instrumented (see \texttt{docs/RAINBOW-RUN-ANALYTICS-BRIEF.md}) --- then every scenario here becomes
a real chart.}

\end{document}
